% Drop-in technical section for the Tessaris/AION system record.
% Implementation state verified on 23 September 2026.
% Compatible with the handover preamble packages hyperref, enumitem and longtable.

\section{Shared Work, Scheduling and Field Delivery System}
\label{sec:tessaris-work-schedule-field-delivery}

\subsection{Purpose and product boundary}

The Work \& Schedule subsystem provides one operational record from an inbound
customer request through scheduling, assignment, field or practitioner delivery,
evidence capture, customer sign-off, quotation and invoicing.  It is intended to
support service businesses whose work may be performed at a customer site, in a
clinic or practice, by video, or from a desktop.  Examples include trades,
consultancies, field engineering, sales appointments, dental or clinical
appointments, and other capacity-constrained professional services.

The central architectural decision is that desktop scheduling and mobile field
work do not maintain separate job databases.  Both clients read and mutate the
same workspace-scoped canonical record through the backend.  The desktop offers
the wider scheduling, intake, stock and customer-history surfaces; the mobile
application exposes the smaller execution surface needed by the assigned person.

Payments are outside the implemented boundary.  The system can retain quotation
and invoice drafts and can describe a future payment boundary, but it does not
charge a card, transmit a payment link, or claim that money moved.  Likewise,
reminders remain prepared, approval-controlled communication records until a
real email or SMS adapter confirms delivery.

\subsection{Implemented components and source files}

\begin{longtable}{p{.28\linewidth}p{.66\linewidth}}
\textbf{Component} & \textbf{Responsibility and source}\ \hline
Canonical domain service &
\path{backend/modules/aion_business/runtime/work_schedule_service.py} owns the
workspace model, normalisation, revisioning, actions, recurrence, conflict
detection, evidence files, generated documents, customer portal, reminders,
customer history and stock movements.\\
Desktop and full-server API &
\path{backend/modules/aion_business/api/work_schedule_api.py} exposes the
FastAPI contract.  Its router is registered by both
\path{backend/desktop_app.py} and \path{backend/main.py}.\\
Desktop workspace &
\path{desktop/mac/src/aion_work_schedule_workspace.js} implements the calendar,
job board, field workspace, booking inbox, customers, reminders, stock and job
detail interactions.  Presentation is in
\path{desktop/mac/src/aion_work_schedule_workspace.css}.\\
Desktop registration &
\path{desktop/mac/src/index.html} loads the workspace assets and
\path{desktop/mac/src/app.js} selects
\texttt{AionWorkScheduleWorkspace.render()} for the sidebar destination.\\
Mobile workspace &
\path{native/ios/PilotApp/Sources/PilotWorkScheduleView.swift} implements the
seven-day schedule, open-work list, job creation, job detail, progress states,
assignment, notes, evidence photographs, checklists, materials, time entries
and directions.\\
Signed mobile client &
\path{native/ios/PilotApp/Sources/PilotNodeClient.swift} sends signed snapshot
and action requests.  \path{PilotPulseView.swift} exposes Work \& Schedule as a
separate mobile application section.\\
Mobile-to-domain bridge &
\path{backend/modules/pilot_unified/workspace_gateway.py} resolves the signed
organisation membership, returns the canonical snapshot, limits admitted action
fields, forwards mutations to \texttt{WorkScheduleService}, and routes real
evidence uploads through the same file service.\\
Tests &
\path{backend/tests/aion_business/test_work_schedule_service.py} covers shared
desktop/mobile jobs, evidence files, PDF generation, clashes, recurrence,
natural-language drafts, revisions, delivery data, commercial numbering,
customer portal actions and reminder preparation.  The signed mobile workspace
contract is additionally covered by
\path{backend/tests/aion_fabric/test_pilot_unified_workspace_gateway.py}.\\
\end{longtable}

\subsection{Canonical storage and data model}

For business workspace $W$, the canonical JSON record is stored at:
\begin{center}\small
\path{<business-containers>/<W>/work_schedule.json}
\end{center}
Uploaded and generated evidence is stored beneath:
\begin{center}\small
\path{<business-containers>/<W>/work_schedule_files/<work-item-id>/}
\end{center}
The current schema identifier is
\texttt{aion.work-schedule.v3}.  A workspace contains:
\begin{itemize}[leftmargin=*]
  \item canonical \texttt{items} and a derived mobile \texttt{jobs} projection;
  \item inbound \texttt{intake} records from website, agent, phone, email or
        customer-portal channels;
  \item people, room, vehicle or equipment \texttt{resources};
  \item consolidated \texttt{customers} and their linked work-item histories;
  \item approval-controlled \texttt{reminders};
  \item an append-style \texttt{stock\_movements} ledger;
  \item hashed, expiring \texttt{portal\_access} grants; and
  \item separate monotonic counters for quote and invoice numbering.
\end{itemize}

Each work item contains identity, title, customer and optional contact fields,
type, state, date, time, duration, assignees, constrained resources, required
skills, location, source and instructions.  Delivery subrecords include work
updates, checklist rows, materials, time entries, dependencies, recurrence,
documents, customer requests, quotation, invoice and customer-signature data.

The service writes JSON through a sibling temporary file followed by replacement,
preventing a partially written canonical file from becoming the visible record.
Every successful mutation advances the workspace revision.  The service accepts
an expected revision and rejects a mismatch with
\texttt{work\_schedule\_revision\_conflict}.  Phone actions and service actions
load and save against the current revision.  The desktop whole-model synchroniser
does not yet supply its last-read revision on every PUT; completing that client
precondition is a remaining multi-writer hardening task.

\subsection{Shared desktop and mobile data flow}

The shared path is:
\begin{enumerate}[leftmargin=*]
  \item the desktop obtains the canonical model from
        \path{/api/aion/business/work-schedule/<workspace>};
  \item an older local-browser record is used only as a one-time migration source
        when the backend workspace is empty;
  \item subsequent desktop changes are persisted to the backend rather than
        treated as an independent source of truth;
  \item the paired phone derives the business workspace from its signed active
        membership and requests \path{/v1/workspaces/work-schedule/snapshot};
  \item the unified gateway returns the service's derived \texttt{jobs}
        projection; and
  \item mobile actions return an updated schedule, so the next desktop or mobile
        read sees the same work, evidence and delivery records.
\end{enumerate}

The mobile client does not accept an arbitrary business identifier from the
interface.  \texttt{PilotNodeClient} signs the request using the possession-bound
phone flow and supplies the membership identifier.  The gateway resolves the
workspace, applies a narrow field allow-list and a 64~KB action-size boundary,
then invokes the canonical service using the authenticated actor identifier.
Evidence uploads are separately bounded to 20~MB before base64 decoding and
durable storage.

\subsection{HTTP API surface}

The desktop and customer-facing API is rooted at:
\begin{center}\small
\path{/api/aion/business/work-schedule}
\end{center}
The implemented routes are:
\begin{longtable}{p{.18\linewidth}p{.40\linewidth}p{.34\linewidth}}
\textbf{Method} & \textbf{Route} & \textbf{Behaviour}\ \hline
GET & \path{/<workspace>} & Read the complete authorised workspace model.\\
PUT & \path{/<workspace>} & Save a normalised model; optionally enforce an expected revision.\\
POST & \path{/<workspace>/actions} & Run a bounded domain action such as create, update, checklist, material, time, stock or reminder preparation.\\
POST & \path{/<workspace>/intake} & Record a reviewable inbound booking request.\\
POST & \path{/<workspace>/items/<id>/files} & Store a real bounded file and attach its metadata to the work item.\\
GET & \path{/<workspace>/items/<id>/files/<document>} & Return an authorised stored evidence or generated document.\\
POST & \path{/<workspace>/items/<id>/documents} & Generate a service report, completion certificate, quote or invoice PDF.\\
POST & \path{/<workspace>/interpret} & Convert supported natural language into a reviewable create or update intention.\\
POST & \path{/<workspace>/items/<id>/portal-access} & Issue a random, expiring customer portal token and retain only its SHA-256 hash.\\
GET/POST & \path{/portal/<token>} and \path{/portal/<token>/actions} & Read the bounded customer projection or record a customer request.\\
GET & \path{/portal/<token>/page} & Render the responsive customer self-service page.\\
POST & \path{/public/<workspace>/booking} & Accept an inbound booking without exposing the schedule.\\
\end{longtable}

Public booking returns only a booking reference and review state.  It does not
return internal work, customers, availability or staff.  Production exposure
still requires the deployment's normal rate limiting, abuse protection and CSRF
policy around this route.

\subsection{Desktop scheduling and back-office workflow}

The desktop workspace includes the following views:
\begin{itemize}[leftmargin=*]
  \item a seven-day calendar with direct creation and drag-to-day rescheduling;
  \item a job board covering new, quotation, approval, scheduled, active and
        completed states;
  \item a field/practitioner view for executing the current workday;
  \item a booking inbox for website, agent and staff-created requests;
  \item customer history consolidated from canonical work items;
  \item a reminder queue that distinguishes prepared communication from sent
        communication; and
  \item a materials ledger for receipt, allocation, use, adjustment and return.
\end{itemize}

A work-detail drawer exposes overview, delivery, updates, files, sign-off and
commercial tabs.  Delivery includes checklist completion, required or used
materials and time entries.  Location-bearing work provides a user-initiated
route link.  The current route is a destination handoff to Google Maps; it is
not a multi-stop optimiser, vehicle tracker or live-traffic scheduler.

The field state machine supports scheduled or new work, on-way/en-route,
in-progress, paused, completed, invoiced, paid and cancelled states.  Buttons
for On my way, Start, Pause, Resume and Complete update the same canonical item
and add an explanatory work-history entry.

\subsection{Availability, recurrence and resources}

Work can carry an assignee, multiple assignee identifiers, resource identifiers
and required skills.  A resource may represent a person, room, vehicle or item
of equipment and records capacity, skills and active state.  Conflict detection
calculates each scheduled interval and reports an overlap when two unfinished
items share an assignee or resource.  The current check detects direct interval
overlap; it does not yet calculate travel time, working hours, leave, skill
substitution or capacity greater than one.

Recurring work supports bounded daily, weekly and monthly series, capped at 52
occurrences.  Each generated item receives its own stable identifier and shares
a series identifier with its origin.  Month generation limits the day to the
valid final day of the target month.  Dependencies can be retained on work
items, but no general critical-path or automatic dependency scheduler is yet
implemented.

\subsection{Work Pilot interpretation}

The Work Pilot bar accepts ordinary wording for booking and changing work.  A
deterministic interpreter is deliberately retained so common operations remain
available when the selected model is offline.  It recognises work types,
customer, assignee, location, duration, common clock forms, relative weekdays,
today, tomorrow and day/month/year dates.  It can prepare new jobs, appointments,
consultations, surveys, visits and meetings.  It also recognises supported
reschedule, move, rebook, assign, cancel, complete and update requests.

An update is applied only when the desktop can resolve exactly one existing work
item from its title, customer or identifier.  No match produces an explicit
failure and multiple matches require the person to choose the record.  The
interpreter prepares structured fields rather than inventing absent facts.

This is a reliable bounded fallback, not broad model-level interpretation.  A
future integration may ask the Vault-selected local or hosted model to propose
the same typed intention, but AION must still resolve the entity, validate the
fields and perform the domain action.  A model must never write the schedule or
claim that an external communication was sent directly.

\subsection{Mobile field and practitioner experience}

The iOS Work section is intentionally narrower than the desktop office surface.
It provides:
\begin{itemize}[leftmargin=*]
  \item a seven-day strip, selected-day bookings, today's count, active count and
        unfinished-work count;
  \item creation of a job or appointment with title, customer, location,
        assignee and start time;
  \item customer, location, appointment time, status and instructions;
  \item on-way, start and complete controls plus editable notes and assignment;
  \item real photo-evidence upload using \texttt{PhotosPicker};
  \item shared checklist completion and addition;
  \item materials read/add controls and shared time-entry recording; and
  \item a user-initiated directions link for work with a recorded location.
\end{itemize}

The iPhone build was compiled using automatic Apple development signing,
installed on the connected physical device and launched under bundle identifier
\texttt{ai.tessaris.pilot}.  Trusting the development certificate and enabling
Developer Mode were device prerequisites for the local development build; they
are not the intended production distribution mechanism.

\subsection{Evidence, reports and customer sign-off}

Desktop and mobile uploads now store bytes rather than filenames or synthetic
evidence counters.  Filenames are reduced to safe basenames, stored names gain a
random prefix and file metadata records type, size, creation time, creator and
evidence kind.  File download resolves the document through its owning work item
before opening the stored path.

The desktop signature pad records the customer name, drawn signature and time on
the work item.  It is a completion or attendance acknowledgement and is not
represented as a substitute for regulated contractual acceptance, identity
verification or qualified electronic signature requirements.

Service reports and completion certificates are generated as PDFs and attached
to the canonical item.  Quotes use identifiers of the form
\texttt{Q-YYYY-NNNN}; invoices use \texttt{INV-YYYY-NNNN}.  Line items are
normalised to description, quantity and unit price.  Subtotal, declared tax rate,
tax and total are calculated using decimal arithmetic before PDF generation.
These are working numbered business documents, but they are not yet posted into
the Finance ledger, reconciled with an accounting provider or qualified for
country-specific tax and invoice law.

\subsection{Customer portal and controlled customer requests}

An authorised desktop user can create a 30-day customer link for one work item.
The service generates a high-entropy token, returns it once and stores only its
SHA-256 digest with the item, customer, creator, creation time and expiry.  Portal
resolution scans for the digest, rejects a revoked or expired grant and returns
only a bounded projection of that work item.  It never returns the wider
schedule, other customers, stock, resources or internal notes.

The responsive portal page shows the work title, customer, status, date and time,
assigned person, location and permitted commercial summary.  The customer can:
\begin{itemize}[leftmargin=*]
  \item approve or decline the current quote;
  \item request another date and time; or
  \item request cancellation and supply a note.
\end{itemize}
Rescheduling and cancellation are stored as
\texttt{needs\_review} customer requests rather than silently rewriting or
deleting office commitments.  Quote approval and decline are recorded on the
commercial record and in the work-item update history.

The portal currently runs from the local Tessaris HTTP service.  A link will be
usable by a remote customer only after an authorised HTTPS deployment, public
routing, abuse controls and privacy review are configured.

\subsection{Reminder and notification boundary}

The desktop scheduler calls \texttt{prepare\_due\_reminders()} for each business
workspace.  The service prepares a 24-hour reminder for work within the next 48
hours and a two-hour reminder for work within two hours, deduplicated by work
item and reminder kind.  It selects email when an email address exists, SMS when
only a phone number exists and otherwise records a manual action.

Email and SMS reminders are stored with
\texttt{status=needs\_approval}; reminders without a deliverable address are
\texttt{manual\_action\_required}.  The approval endpoint
\path{/<workspace>/reminders/<reminder>/approve-and-send} binds a named approver
to the prepared reminder.  Email is submitted through the existing Resend
adapter and SMS through Twilio Messaging.  A successful submission must return
a provider request identifier before the reminder becomes
\texttt{submitted\_to\_provider}; absent credentials, provider failure or a
missing receipt records \texttt{failed}.  Reusing a reminder after it leaves
\texttt{needs\_approval} is rejected, preventing an accidental duplicate send.
Submission is not described as delivery: provider delivery-event reconciliation
remains a separate production step.

The Vault now exposes a purpose-separated integration catalogue.  Gmail and
Microsoft Outlook are mailbox providers; Mailchimp Marketing manages audiences
and campaigns; Resend is the recommended operational email transport; Twilio
Messaging provides SMS; Google Maps Routes is the preferred routing candidate;
Xero and QuickBooks Online are ledger candidates; and HubSpot or Pipedrive are
the primary CRM candidates.  Catalogue, adapter, configured, connected and
live-verified are different states.  Only Gmail currently has the established
local OAuth health bridge.  Outlook has a bounded Microsoft Graph draft/send
adapter, while OAuth registration and token refresh still have to be connected
through Vault before use.

\subsection{Stock, materials and time}

Per-job materials record name, quantity, unit, optional unit cost and whether the
material was used.  The workspace stock ledger records receive and return as
positive quantities and allocate, use or downward adjustment as negative
quantities.  Every movement has a stable identifier, optional work-item link,
note, actor and time.  The desktop derives the displayed balance from these
movements rather than mutating a separate total.

Time entries record bounded minutes, purpose, actor and timestamp against the
work item.  The delivery view sums those entries for field and back-office use.
This supports operational time capture but is not a payroll timesheet, wage
calculation or labour-law compliance system.  Stock similarly remains an
operational ledger; reorder rules, suppliers, purchase orders, serial or batch
tracking, returns-to-supplier and accounting valuation remain future modules.

\subsection{Provider-neutral accountant harness}

The Finance integration has been extended beyond a list of accounting brands.
\path{accounting_provider_contract.py} defines one provider boundary and a
canonical set of accounts, invoices and supplier bills, payments, bank
transactions, journals, tracking categories and projects.  Provider adapters
may retrieve and translate facts; they do not decide tax treatment, select an
ambiguous nominal account or execute an unapproved posting.

\path{quickbooks_online_accounting_adapter.py} reads QuickBooks Online accounts,
invoices, bills, payments, purchases and journal entries through the Intuit
Accounting API and maps them to the canonical ledger.  It preserves the Intuit
company realm, provider identifiers and source collection in every source
reference.  \path{sage_accounting_adapter.py} performs the equivalent mapping
for Sage Business Cloud Accounting ledger accounts, sales and purchase
invoices, contact payments, bank-account transactions and ledger entries.
\path{freeagent_accounting_adapter.py} reads categories, contacts, invoices,
bills, bank accounts, bank transactions and journal sets, including the
explained/unexplained evidence attached to FreeAgent bank activity.
\path{freshbooks_accounting_adapter.py} reads service-business invoices,
supplier bills and their payment evidence, customer payments and expenses and,
when the OAuth identity supplies the separate
business UUID, the FreshBooks chart of accounts.  The Vault connection must
retain both the FreshBooks account identifier and business UUID; a missing UUID
is reported as missing account evidence rather than replaced with an invented
chart.  \path{zoho_books_accounting_adapter.py} reads the chart of accounts,
sales invoices, bills, customer and vendor payments, bank transactions and
journals.  It retains the organisation identifier and OAuth-provided regional
API domain so European and other non-US Zoho tenants are not routed to the
wrong data centre.
Access and refresh credentials are not accepted by the normalisation boundary
and must remain in Vault-managed provider connection state.

\path{accountant_harness_service.py} writes a read-only, hash-bound sync receipt,
raw evidence hash and canonical Finance ledger.  That ledger is deliberately
compatible with the existing Xero Finance matching engine, so all six
providers share deterministic amount, date, reference and counterparty
matching, duplicate detection, exceptions and human review.  No provider write
is performed by ingestion or reconciliation.  Future mutations must retain the
existing Xero control sequence: prepare an exact payload, bind approval to its
hash, use an idempotency key, execute through the selected provider adapter,
read the created or changed object back independently, and bind the receipt to
the approved hash.

The Finance integration API exposes the accounting-provider capability catalog,
an internal post-authorisation snapshot-ingestion boundary and a provider-neutral
reconciliation route.  QuickBooks Online, Sage, FreeAgent, FreshBooks and Zoho
Books are therefore
\emph{adapter-ready}, not connected or live-verified.  A live claim requires
provider application registration, OAuth connection through Vault, sandbox
acceptance, an approved write test, read-back reconciliation, failure testing
and credential revocation testing.  Provider payment records may be imported
as accounting evidence for reconciliation; this does not enable Tessaris to
charge a card, initiate a bank payment or otherwise execute payment processing.

\subsection{Safety and authority properties}

The implementation enforces the following distinctions:
\begin{itemize}[leftmargin=*]
  \item a booking request is not a scheduled commitment until an authorised
        action creates or updates a work item;
  \item a prepared reminder is not a sent notification;
  \item a customer reschedule or cancellation is a request awaiting review;
  \item a generated quote or invoice is not an accounting-ledger posting;
  \item a recorded payment outcome, where legacy UI data exists, is not evidence
        that Tessaris charged a payment instrument; and
  \item model interpretation proposes typed fields while the governed service
        performs the mutation.
\end{itemize}

Mobile mutations are bound to a signed membership-derived workspace.  Public
booking cannot read the schedule.  Portal tokens are random, hashed at rest,
scoped to one work item and time-limited.  Evidence has a size boundary and safe
stored name.  Domain operations reject missing work items, unsupported actions,
invalid money values, invalid stock operations and duplicate work identifiers.

\subsection{Verification and installed state}

The implementation was verified through Python compilation, JavaScript syntax
checking, the Work Schedule domain suite and unified mobile gateway suite.  At
the recorded closure, 41 relevant tests passed.  The iOS target compiled and
signed successfully, and the resulting application was installed and launched
on the connected iPhone.  The Electron desktop package was rebuilt, installed
at \path{/Applications/Tessaris.app}, launched and checked against the live
desktop backend.  The installed interface exposed the shared-mobile indicator,
calendar, job board, field workspace, booking inbox, customers, reminders and
stock/materials views.

Automated verification covers structure and bounded domain behaviour.  It does
not qualify remote customer hosting, real message delivery, route optimisation,
tax compliance, accounting posting, payment processing or industry regulation.

\subsection{Remaining production and expansion work}

The following are deliberately not closed by this implementation:
\begin{enumerate}[leftmargin=*]
  \item connect production Resend and Twilio credentials, verify sending domains
        and numbers, and reconcile delivered, bounced, failed and opted-out
        provider events after submission;
  \item deploy the customer portal behind authorised public HTTPS routing,
        authentication or equivalent possession controls, rate limiting and
        privacy policy;
  \item add map-native multi-stop optimisation, travel-time constraints, live
        traffic and vehicle or practitioner routing;
  \item complete bidirectional linkage to the Sales Revenue Spine and Finance
        customer/invoice ledgers rather than relying only on derived work history;
  \item supply line-item editing, credit notes, tax jurisdiction, due dates,
        accounting export, ledger posting and reconciliation;
  \item add purchasing, suppliers, reorder thresholds, stock locations, batches,
        serials, returns and vertical-specific inventory controls;
  \item add richer checklist templates, measurement schemas, timer start/stop,
        labour approval and practitioner-specific clinical or regulated records;
  \item route broad language through the Vault-selected accepted model while
        retaining deterministic parsing, entity resolution and governed domain
        execution as the fail-safe path;
  \item make the desktop synchroniser send an expected revision on every write
        and provide an explicit merge experience after concurrent edits; and
  \item connect a certified payment provider only through a separately approved,
        receipt-backed payment capability.
\end{enumerate}

The correct production claim is therefore that Tessaris now has a functional,
shared desktop/mobile work-and-delivery foundation with real local persistence,
evidence, controlled customer self-service and operational records.  External
communications, public deployment, accounting qualification, route optimisation,
specialist vertical workflows and payments remain integration or qualification
work and must not be inferred from the presence of their internal records.
